% ============================================================================
\section{TSS 初始化}
\label{sec:tss-init}
% ============================================================================

\texttt{tss\_init()} 做 4 件事：清零、设置 SS0:ESP0、写 GDT 描述符、加载 TR。

\codefile{src/kernel/arch/tss.c}
\begin{lstlisting}[style=cstyle]
static tss_t g_tss;
extern uint32_t stack_top;  /* boot.asm 中的内核栈顶 */

void tss_init(void) {
    /* 1. 清零 TSS（所有不用的字段 = 0） */
    kmemset(&g_tss, 0, sizeof(g_tss));

    /* 2. 设置 Ring 0 栈段选择子 */
    g_tss.ss0 = GDT_KERNEL_DATA_SEL;  /* 0x10 */

    /* 3. 设置 Ring 0 栈顶指针 */
    g_tss.esp0 = (uint32_t)(uintptr_t)&stack_top;

    /* 4. I/O 位图偏移 = sizeof(tss)，表示没有 I/O 位图
     *    → Ring 3 执行 IN/OUT 指令会触发 #GP */
    g_tss.iomap_base = sizeof(tss_t);

    /* 5. 在 GDT[5] 写入 TSS 描述符 */
    gdt_set_tss_entry((uint32_t)(uintptr_t)&g_tss,
                       sizeof(tss_t) - 1);  /* limit = 103 */

    /* 6. ltr 加载到 TR 寄存器 */
    ltr(GDT_TSS_SEL);  /* 0x28 */
}
\end{lstlisting}

% ----------------------------------------------------------------------------
\subsection{\texttt{ltr} 指令封装}
% ----------------------------------------------------------------------------

\begin{lstlisting}[style=cstyle]
static inline void ltr(uint16_t sel) {
    __asm__ volatile("ltr %0" : : "r"(sel));
}
\end{lstlisting}

\begin{cpustate}
\texttt{ltr 0x28} 执行后：
\begin{itemize}
  \item TR 寄存器指向 GDT[5] 的 TSS 描述符
  \item CPU 内部缓存 TSS 的基址和 limit
  \item GDT[5] 的 type 自动从 \texttt{0x89}（Available）变为 \texttt{0x8B}（Busy）
  \item 之后发生 Ring 3 → Ring 0 切换时，CPU 会从 \texttt{g\_tss.esp0} 获取内核栈
\end{itemize}
\end{cpustate}

% ----------------------------------------------------------------------------
\subsection{\texttt{tss\_set\_kernel\_stack()} —— 为进程切换预留}
% ----------------------------------------------------------------------------

将来做进程切换时，每个进程有独立的内核栈。切换进程后需要更新 TSS.ESP0：

\begin{lstlisting}[style=cstyle]
void tss_set_kernel_stack(uint32_t esp0) {
    g_tss.esp0 = esp0;
}
\end{lstlisting}

当前阶段只有一个内核栈，此函数暂时不会被调用，但接口先准备好。

% ----------------------------------------------------------------------------
\subsection{初始化时机与顺序}
% ----------------------------------------------------------------------------

\codefile{src/kernel/core/kmain.c}
\begin{lstlisting}[style=cstyle]
gdt_init();     /* GDT 先就位 */
idt_init();     /* IDT 先就位（处理 Ring 3 异常） */
tss_init();     /* TSS 依赖 GDT，且需要 IDT 已安装 */
\end{lstlisting}

\begin{whybox}
为什么 \texttt{tss\_init()} 必须在 \texttt{gdt\_init()} 之后？

因为 \texttt{tss\_init()} 会调用 \texttt{gdt\_set\_tss\_entry()} 修改 GDT[5]，
然后用 \texttt{ltr} 加载 TSS——\texttt{ltr} 会去 GDT 查找描述符。
如果 GDT 尚未初始化（\texttt{gdtr} 没有指向有效表），\texttt{ltr} 会触发异常。
\end{whybox}
